% Main-text fragment; supplementary proofs are in supplement_theory.tex.
% Requires amsmath, amssymb, booktabs and natbib.
\section{Population geometry and the information lost in a directional representation}
\label{sec:origin-theory}

The geometric cases introduced in Fig.~\ref{fig:geometry} are formalized by two continuous half-annular class supports. The displacement $h$ changes the position of the second support; the normalization origin $c$ is a separate choice that determines which observations share a direction. These definitions make it possible to distinguish ambiguity in the original data from the additional Bayes error created by the representation.

\begin{samepage}
Let $a=0.8$, $b=1.2$, $m_0=(0,0)$ and $m_1=(1,h)$, where $h$ is the upward center displacement. The open class supports are
\begin{align}
 S_0&=\{x:a<\|x\|<b,\ x_2>0\},\nonumber\\
 S_1&=\{x:a<\|x-m_1\|<b,\ x_2<h\}.
 \label{eq:halfmoon-supports}
\end{align}
\end{samepage}
Each support has area $H=\pi(b^2-a^2)/2=2\pi/5$. The primary generator draws points uniformly over the area of each band, so its density is $f_i^A=\mathbf1_{S_i}/H$. A second law draws radius and angle uniformly, giving $f_i^R(x)=\mathbf1_{S_i}(x)/[\pi(b-a)\|x-m_i\|]$; it tests whether the geometric conclusions survive a different allocation of probability within the same supports. Neural experiments use equal priors and area-uniform sampling. The geometric positivity results apply to any positive priors $\pi_i$ and densities positive almost everywhere on the respective supports.

The two-point collision in the Introduction demonstrates that normalization is not one-to-one. It does not by itself prove a positive population error, because a point or single ray has zero probability under the continuous sampling laws. The relevant condition is conflict between Bayes decisions across a set of directions with positive angular measure. The following identity measures that conflict from the class densities.

\paragraph{Why one ray can force classification error.}
Direction-only classification must use the same decision for every radius associated with an angle. The class-weighted preference is expressed by the signed density $d=\pi_0f_0-\pi_1f_1$: its positive and negative regions favor classes 0 and 1, respectively. Let $u_\theta=(\cos\theta,\sin\theta)$ and $[z]_+=\max(z,0)$. Integration along the ray from $c$ uses the polar area factor $r$, because an angular slice contains more area at a larger radius. The opposing class-preference masses per unit angle are
\begin{equation}
\begin{aligned}
 A_c(\theta)&=\int_0^\infty[d(c+ru_\theta)]_+r\,dr,\\
 B_c(\theta)&=\int_0^\infty[-d(c+ru_\theta)]_+r\,dr.
\end{aligned}
\label{eq:ray-masses-main}
\end{equation}
Let $R_X^*$ and $R^*_{T_c}$ denote the minimum classification risks from raw coordinates and directions, respectively. Using the masses in~\eqref{eq:ray-masses-main}, the binary Bayes-risk identity yields
\begin{equation}
 \Delta(c):=R^*_{T_c}-R_X^*
 =\int_0^{2\pi}\min\{A_c(\theta),B_c(\theta)\}\,d\theta.
 \label{eq:ray-loss-main}
\end{equation}
The functions $A_c(\theta)$ and $B_c(\theta)$ are signed decision-mass \emph{densities per unit angle}, not probabilities on individual zero-width rays. Their minimum describes opposing preference cancelled at each angle; only integration across angles gives the excess Bayes risk $\Delta(c)$, a probability. Positive loss therefore requires conflicting preferences on a set of directions of positive angular measure. This specializes decision-preservation theory \citep{Silva2012,Silva2024} to a moving normalization origin; Supplementary Section~\ref{supp:bayes} derives the identity.

\paragraph{Theorem 1: loss at every finite origin.}
A collision seen from one origin might disappear when the origin moves. The stronger question is whether some viewpoint separates the two half-rings into different sets of directions. Their geometry answers that question for the positive displacements below: despite disjoint original class regions, every finite viewpoint merges conflicting decisions. For $0<h\leq\sqrt{3/5}$, the supports have disjoint interiors and
\begin{equation}
 R_X^*=0,\qquad \Delta(c)>0
 \quad\text{for every }c\in\mathbb R^2.
 \label{eq:all-origin-disjoint}
\end{equation}
At $h=0$, every origin with $c_2=0$ preserves zero risk. More generally, for $h<0$, origins in the separating strip $h\leq c_2\leq0$ preserve zero risk. At $h=1$, the supports overlap and
\begin{equation}
 R_X^*>0,\qquad \Delta(c)>0
 \quad\text{for every }c\in\mathbb R^2.
 \label{eq:all-origin-overlap}
\end{equation}

\emph{Proof idea.}
Horizontal slicing shows that the disjoint supports satisfy $\sqrt{a^2-h^2}\geq1/5$ in the stated range. Yet each connected half-ring intrudes into the other one's convex hull. Were their directional images disjoint from some origin, a separating angular cone would contain one region's convex hull while excluding the other, contradicting this intrusion. This adapts the geometry of Hurtado et al.~\citep{Hurtado2001}; positive class densities convert the forced angular overlap into positive risk in~\eqref{eq:ray-loss-main}. At $h=1$, the same argument uses connected class-exclusive regions to establish additional loss despite the overlap already present in raw coordinates. Supplementary Sections~\ref{supp:geometry}--\ref{supp:overlap} provide the explicit geometric witnesses, inequalities and complete proofs.

% Suggested main-text replacement/extension for the old positive-infimum theorem.
% The supplement fragment defines supp:projection-boundary.
Finite-origin positivity does not imply a globally positive lower bound. A sequence of strictly positive values may converge to zero along increasingly distant centers. This distinction is fundamental when the goal is to establish an irreducible error for \emph{every} possible normalization origin. Viewed from afar, rays through a bounded population become almost parallel, so their retained differences are associated with projections perpendicular to the viewing direction. The statistical statement needed here concerns the optimal classification risks of those representations rather than a visual resemblance between rays. The uniform theorem below controls all viewing directions simultaneously and completes the origin-search problem at infinity.

\paragraph{Uniform projection limit at distant origins.}
Fix a reference point \(p\), a unit viewing vector \(e\), and the orthogonal projection \(P_e\) onto the hyperplane perpendicular to \(e\). For binary populations with densities in dimension \(n\geq2\), the optimal risk of direction normalization about \(p+Re\) approaches the risk of \(P_e\) uniformly over \(e\) as \(R\) increases:
\begin{equation}
 \lim_{R\to\infty}\sup_e
 \left|R^*_{T_{p+Re}}-R^*_{P_e}\right|=0.
 \label{eq:projection-limit-main}
\end{equation}
Write $\Delta_\infty(e)=R^*_{P_e}-R_X^*$. The loss surface, completed by these projection limits, is continuous on a compact domain. Consequently,
\begin{equation}
 \inf_c\Delta(c)>0
 \quad\Longleftrightarrow\quad
 \begin{cases}
 \Delta(c)>0&\text{all finite }c,\\
 \Delta_\infty(e)>0&\text{all directions }e.
 \end{cases}
 \label{eq:global-criterion-main}
\end{equation}
Thus optimizing the origin cannot make loss vanish when both finite radial views and all limiting linear projections lose decisions.

\emph{Half-moon consequence.}
The connected winner regions in Theorem~1 have overlapping projection intervals in every direction. Equations~\eqref{eq:projection-limit-main}--\eqref{eq:global-criterion-main} therefore strengthen its conclusion to
\begin{equation}
 \inf_{c\in\mathbb R^2}\Delta(c)>0
 \quad\text{for }0<h\leq\sqrt{3/5}\text{ and for }h=1.
 \label{eq:global-positive}
\end{equation}
This holds for every fixed pair of positive priors and densities positive almost everywhere on their respective supports. The proof controls signed probability mass in total variation, so it establishes convergence of the optimal decision risk. Supplementary Section~\ref{supp:projection-boundary} proves the general result and gives an explicit convergence bound for bounded-support densities of bounded variation.

% Standalone main-text replacement paragraph. Definitions of d and Delta precede it.
\paragraph{A numerical lower certificate for the whole plane.}
A certificate converts the qualitative positivity theorem into a numerical lower bound valid for every origin. Such a bound can be constructed from two interior regions where one class is favored by a fixed probability margin. Their relative positions make conflicting rays unavoidable, while their sizes bound the amount of probability merged on those rays. In the plane, let $K_+,K_-$ be compact connected cores whose radius-$\rho$ neighborhoods satisfy $d\geq\kappa_+>0$ and $d\leq-\kappa_-<0$, respectively. Suppose each core meets the other's convex hull, their convex hulls contain a common disk $B(p,\delta)$, and both cores lie within distance $M$ of $p$. Writing $\kappa=\min(\kappa_+,\kappa_-)$ gives
\begin{equation}
 \inf_c\Delta(c)\geq\kappa\min\left\{
 \frac{\sqrt3\rho^3}{6M},\;2\rho\delta\right\}>0.
 \label{eq:constructive-global}
\end{equation}
For bounded origins, mutual convex-hull intrusion forces rays through both cores; their neighborhoods provide an angular interval with conflicting class mass. For distant origins, the shared convex-hull disk supplies the bound. Thus the certificate does not rest on a finite search grid. Supplementary Section~\ref{supp:constructive-global} proves the result and optimizes the split between these regions. For equal-prior area sampling, the conservative whole-plane certificates exceed $0.0205$ percentage points (pp) at $h=0.5$ and $0.00119$~pp at $h=1$. They bound $\inf_c\Delta(c)$, the best loss over \emph{all} origins; by comparison, the midpoint loss at $h=0.5$ is $10.7479\%$ for \emph{one fixed} origin (Table~\ref{tab:neural-main}). Eroded half-annuli supply the required cores under both sampling laws and any fixed positive priors.

\paragraph{Computing the best possible error at a chosen origin.}
An ideal directional classifier needs no neural training: for each viewing angle, one can integrate the class probabilities carried by that direction and choose the more probable label. For the half-annuli, circle tangencies and boundary intersections partition directions into finitely many cells. Radial support order is fixed inside a cell, so one representative ray determines whether opposite preferences coexist. Exact radial primitives then reduce Bayes risk to one angular integral. Ball arithmetic and a bounded-variation truncation bound certify its value. Supplementary Sections~\ref{supp:finite}--\ref{supp:bv} give the finite preservation test, integration formulas and error bounds. At equal priors, the density ratio $f_i^R/f_i^A\in[5/6,5/4]$ also gives $\Delta_R(c)\geq(5/6)\Delta_A(c)$.

\paragraph{Recovering the decision with one radius bit.}
The all-origin loss result rules out fixing the midpoint case by moving the origin, but it leaves open whether a small amount of distance information can repair the representation. The proposed bit answers only whether an observation is near or far from the fixed midpoint; it does not report the full distance. At $h=0.5$, direction alone has positive loss at every origin. At $c=(1/2,1/4)$, append $q=\mathbf1\{\|x-c\|>4/5\}$. For $q=0$, the sign of $9u_1+10u_2$ determines the class; for $q=1$, the sign of $u_2-u_1$ determines it. Positive signs identify class~0. Figure~\ref{fig:bit-repair} shows how the radial bit assigns observations to the two angular rules. Together, the two rules correctly label the complete class supports under their respective bit conditions, giving zero Bayes risk under either sampling law and any positive priors. Supplementary Section~\ref{supp:one-bit} proves the separating inequalities. Since zero additional radial bits leave the positive direction-only floor, one bit is the minimal radial addition at this geometry.

\paragraph{How accurately must radius be measured?}
With the same angular rules, recovery is exact precisely for $\alpha\leq t\leq\beta$, where $\alpha\approx0.7433034374$ and $\beta\approx0.8120914734$ (Threshold Addendum, Proposition~1). If the measured radius is $\widehat r=r+e$, $|e|\leq\varepsilon$, every support label is preserved for every admissible error function exactly when
\begin{equation}
\alpha+\varepsilon\leq t\leq\beta-\varepsilon.
\label{eq:radius-error}
\end{equation}
At the experimentally fixed $t=0.8$, the decoder tolerates bounded radius error up to $\beta-0.8\approx0.0120914734$ coordinate units. The different threshold $t_\star=(\alpha+\beta)/2\approx0.7776974554$ maximizes that guarantee at $(\beta-\alpha)/2\approx0.0343940180$ units (Addendum, Corollary~1). The neural experiment uses exact-radius bits at $t=0.8$; it does not test noisy radius measurements.
